% Part: sets-functions-relations
% Chapter: sets
% Section: basics

\documentclass[../../../include/open-logic-section]{subfiles}

\begin{document}

\olfileid{sfr}{set}{bas}
\olsection{সদস্যভিত্তিক সমতার নীতি}

কতগুলি বস্তুর সমষ্টিকে {একটি} বস্তু হিসেবে বিবেচনা করলে তাকে \emph{সেট} বলে।
যে বস্তুগুলি নিয়ে সেটটি গঠিত, সেগুলিকে সেটটির \emph{উপাদান} বা
\emph{সদস্য} বলা হয়। $x$ যদি কোনও সেট~$A$-এর !!a{element} হয়, তবে আমরা
$x \in A$ লিখি; তা না হলে $x \notin A$ লিখি। যে সেটের কোনও
!!{element}s নেই, তাকে \emph{শূন্য} সেট বলে এবং
``$\emptyset$'' দিয়ে নির্দেশ করা হয়।

\begin{explain}
আমরা সেটটি কীভাবে \emph{নির্দেশ} করি, তার !!{element}s কী \emph{ক্রমে}
সাজাই, কিংবা তার !!{element}s \emph{কত বার} গুনি, তা বিবেচ্য নয়।
বিবেচ্য কেবল এই যে, তার !!{element}s কী কী। নিচের নীতিতে আমরা
এই কথাটিই সুনির্দিষ্টভাবে প্রকাশ করি।
\end{explain}

\begin{defn}[সদস্যভিত্তিক সমতার নীতি]
  $A$ ও $B$ সেট হলে, $A = B$ হবে যদি এবং কেবল যদি
  $A$-এর প্রতিটি !!{element} $B$-এরও !!a{element} হয় এবং
  বিপরীত দিক থেকেও একই কথা সত্য হয়।
\end{defn}

সদস্যভিত্তিক সমতার নীতি কিছু প্রতীক ব্যবহারের ভিত্তি দেয়। সাধারণভাবে,
কতগুলি বস্তু $a_{1}$, \dots, $a_{n}$ দেওয়া থাকলে,
$\{a_{1}, \dots, a_{n}\}$ হল \emph{সেই নির্দিষ্ট} সেট, যার
!!{element}s হল $a_1, \ldots, a_n$। এখানে ``\emph{সেই নির্দিষ্ট}''
কথাটির উপর জোর দিচ্ছি, কারণ সদস্যভিত্তিক সমতার নীতি অনুযায়ী এমন সেট
\emph{একটিই} হতে পারে। এই নীতি নিচের সমতাও সমর্থন করে:
  \[
    \{a, a, b\} = \{a, b\} = \{b,a\}.
  \]
এতে আগের কথাটিই স্পষ্ট হয়: সেট বিবেচনার সময় তার !!{element}sের
ক্রম কিংবা সেগুলি কত বার উল্লেখ করা হয়েছে, তা নিয়ে আমাদের মাথা ঘামাতে হয় না।

\begin{tagblock}{novice}
\begin{ex}
কতগুলি বস্তু দেওয়া থাকলে, সেগুলিকে একত্র করে একটি সেট গঠন করা যায়।
যেমন, রিচার্ডের ভাইবোনদের সেটে একজন ব্যক্তি আছেন; সেটটিকে আমরা
$S=\{\textrm{রুথ}\}$ লিখতে পারি। $4$-এর চেয়ে ছোট ধনাত্মক পূর্ণসংখ্যাগুলির
সেট হল $\{1, 2, 3\}$; আবার সেটিকে $\{3, 2, 1\}$,
এমনকি $\{1, 2, 1, 2, 3\}$ হিসেবেও লেখা যায়। সদস্যভিত্তিক সমতার নীতি
অনুযায়ী এগুলি একই সেট। কারণ, $\{1, 2, 3\}$-এর প্রতিটি !!{element}
$\{3, 2, 1\}$-এরও (এবং $\{1, 2, 1, 2, 3\}$-এরও) !!a{element},
আর বিপরীত দিক থেকেও একই কথা সত্য।
\end{ex}
\end{tagblock}

প্রায়ই কোনও সেটের !!{element}sের সাধারণ কোনও ধর্মের সাহায্যে সেটটিকে
নির্দেশ করব। তার জন্য এই সংক্ষিপ্ত প্রতীক ব্যবহার করব:
$\Setabs{x}{\phi(x)}$। এখানে $\phi(x)$ হল সেই ধর্ম, যা~$x$-এর থাকতে হবে,
যাতে তাকে সেটটির !!{element}sের মধ্যে গণ্য করা যায়।

\begin{tagblock}{novice}
\begin{ex}
আমাদের উদাহরণে $S$-কে এভাবেও নির্দেশ করা যেত:
\[
S = \Setabs{x}{x \text{ রিচার্ডের ভাই বা বোন}}.
\]
\end{ex}
\end{tagblock}

\begin{tagblock}{math}
\begin{ex}
কোনও সংখ্যাকে \emph{নিখুঁত} বলা হয় যদি এবং কেবল যদি সেটি তার প্রকৃত
ভাজকগুলির যোগফলের সমান হয়। প্রকৃত ভাজক বলতে সেই সংখ্যাগুলিকে বোঝায়,
যেগুলি তাকে নিঃশেষে ভাগ করে, কিন্তু সংখ্যাটির সমান নয়। যেমন, $6$
নিখুঁত, কারণ তার প্রকৃত ভাজকগুলি হল $1$, $2$ এবং~$3$, আর
$6 = 1 + 2 + 3$। বস্তুত, $10$-এর চেয়ে ছোট ধনাত্মক পূর্ণসংখ্যাগুলির
মধ্যে কেবল $6$-ই নিখুঁত। তাই সদস্যভিত্তিক সমতার নীতি ব্যবহার করে বলতে পারি:
\[
	\{6\} = \Setabs{x}{x\text{ নিখুঁত এবং }0 \leq x \leq 10}
\]
ডান দিকের প্রতীকটি আমরা পড়ি: ``সেই সব $x$-এর সেট, যাদের জন্য $x$
নিখুঁত এবং $0 \leq x \leq 10$''। এই সমতা নিশ্চিত করে যে, সেট
বিবেচনার সময় সেটিকে কীভাবে নির্দেশ করা হয়েছে, তা বিবেচ্য নয়।
আরও সাধারণভাবে, সদস্যভিত্তিক সমতার নীতি নিশ্চিত করে যে, $\phi(x)$
ধর্মবিশিষ্ট $x$-গুলির সেট সব সময় একটিই হয়। তাই এই নীতির ভিত্তিতেই
$\Setabs{x}{\phi(x)}$-কে $\phi(x)$ ধর্মবিশিষ্ট $x$-গুলির
\emph{সেই নির্দিষ্ট} সেট বলা যায়।
\end{ex}
\end{tagblock}

সদস্যভিত্তিক সমতার নীতি দুটি সেট অভিন্ন প্রমাণ করার একটি উপায় দেয়:
$A = B$ দেখাতে হলে দেখাও যে, যখনই $x \in A$, তখনই $x \in B$,
এবং যখনই $y \in B$, তখনই $y \in A$।

\begin{prob}
প্রমাণ করো যে, শূন্য সেট একটির বেশি হতে পারে না। অর্থাৎ, দেখাও যে,
$A$ ও $B$ এমন সেট হলে যাদের কোনও !!{element}s নেই, $A = B$ হবে।
\end{prob}

\end{document}
